\begin{abstract}
Text watermarking schemes of the SynthID-Text class embed a statistical signal in the sampling process of a language model, so the mark lives in the choice of tokens rather than in attached metadata. A growing literature measures how well such marks survive individual transformations: paraphrasing, back-translation, token substitution, and regeneration. That literature evaluates one transformation at a time, on its own operating point, and rarely reports how much meaning the transformation costs. This paper asks a joint question. Given a fixed semantic-preservation floor, which single transformation, or which ordered chain of at most three transformations, removes the watermark most reliably? We formalise the problem as a constrained search over an attack catalogue, measure removal against a same-config detector under the generating key and at a corpus-level true-positive rate, and select chains with a document-level composite of BERTScore and COMET, with the other metrics reported beside it. The design evaluates the selection on a held-out split, clusters every interval by prompt, and audits the chains the search skips. We release the catalogue, the search procedure, and the evaluation harness so that the frontier between removal and preservation can be reproduced and extended.
\end{abstract}
